Le but du problème est l’étude de la demande et de l’offre pour un nouveau produit
de grande consommation.

\begin{itemize}[$\bullet$]
	\item $x$ désigne le prix unitaire en euros du produit ;
	\item $f$ désigne la demande (la quantité de produit demandée par les consommateurs),
en milliers d’unités ;
	\item $h$ désigne l’offre (la quantité de produit offerte sur le marché par les producteurs),
en milliers d’unités.
\end{itemize}





On considère l’équation différentielle (E) : $y' + 0,4 y = 0,4x-1$ où $y$ désigne une fonction de la variable réelle $x$ définie et dérivable sur $\R$, $y'$ sa fonction dérivée.

\begin{enumerate}
	\item Donner, puis résoudre l'équation homogène (E$_0$) :\sautp
	
	\Lignes{2}
	\item A l’aide du logiciel Géogébra, déterminer les nombres réels $a$ et $b$ pour que la fonction $p$ définie sur $\R$ par : $p(x) = a~x+b$ soit une solution particulière de l'équation différentielle (E), et donner l’expression de $p$. 

	
	On trouve : \hfill $a=$\makebox[0.1\linewidth]{\dotfill} \hfill $b=$\makebox[0.1\linewidth]{\dotfill} \hfill $p(x)=$\makebox[0.3\linewidth]{\dotfill}

	\smallskip
	
	\appelprof
	
	\item En déduire l'ensemble des solutions de (E) : \sautp
	
	\Lignes{1}
	\item \begin{enumerate}
		\item Déterminer, par un calcul, la solution $f$ de (E) telle que $f(0)=10$ : \sautp
	
	\Lignes{3}
		\item Interpréter, dans le contexte de l'énoncé, la condtion initiale $f(0)=10$ :\sautp
	
	\Lignes{2}
	\end{enumerate}
	
	
	\item L'offre est représentée par la fonction $h(x) = \ell n(3x + 0, 9)$. Le prix de vente est fixé comme la valeur de $x$ pour laquelle l'offre est égale à la demande.\sautm
	
	À l'aide de Géogebra, déterminer le prix de vente du produit : \dotfill
	

\end{enumerate}